%=========== ΛΥΣΗ - 60ο ΘΕΜΑ Δ ===========
\begin{thema}{Δ}
\begin{erwthma}
\item Η $f$ είναι παραγωγίσιμη στο $(0,+\infty)$ και για κάθε $x>0$ είναι
\[
f'(x)=\frac{1-\ln x}{x^2}.
\]
Επειδή $x^2>0$, το πρόσημο της $f'$ είναι το πρόσημο της παράστασης $1-\ln x$. Έχουμε:
\begin{itemize}
\item $f'(x)>0$ για $0<x<e$,
\item $f'(e)=0$,
\item $f'(x)<0$ για $x>e$.
\end{itemize}
Στον παρακάτω πίνακα βλέπουμε το πρόσημο της $f'$ και τη μονοτονία της $f$:
\begin{center}
\begin{tikzpicture}
\tkzTabInit[color,lgt=2,espcl=1.8,colorC = maincolor!50!white,colorL = maincolor!20!white,colorV = maincolor!50!white]%
{$x$ / .8,$f'(x)$ /0.8,$f(x)$/1}%
{$0$,$e$,$+\infty$}%
\tkzTabLine{d, +, z, -, }
\tkzTabVar{D-/{-\infty},+/{\frac{1}{e}},-/{0}}
\end{tikzpicture}
\end{center}

Άρα η $f$ είναι \textbf{γνησίως αύξουσα} στο $(0,e]$ και \textbf{γνησίως φθίνουσα} στο $[e,+\infty)$. Επομένως παρουσιάζει ολικό μέγιστο στο $x_0=e$, με τιμή
\[
{f(e)=\frac{1}{e}}.
\]
Δεν παρουσιάζει ελάχιστο, αφού
\[
\lim_{x\to0^+}f(x)=-\infty.
\]

Συνεπώς, για κάθε $x>0$,
\[
{\frac{\ln x}{x}\leq\frac{1}{e}},
\]
με ισότητα μόνο για $x=e$.

\item Από την προηγούμενη ανισότητα έχουμε
\[
\frac{\ln x}{x}\leq\frac{1}{e}.
\]
Επειδή $ex>0$, προκύπτει
\[
{e\ln x\leq x}
\]
για κάθε $x>0$. Η ισότητα ισχύει μόνο για $x=e$.

Για $x=\pi$ έχουμε $\pi\neq e$, οπότε η ανισότητα είναι αυστηρή:
\[
e\ln\pi<\pi.
\]
Άρα
\begin{align*}
\ln\pi^e&<\ln e^\pi\\
\Leftrightarrow
{\pi^e<e^\pi}.
\end{align*}

\item Με βάση τη μονοτονία και τα όρια της $f$, διακρίνουμε τις εξής περιπτώσεις:
\begin{itemize}
\item Αν $k<0$, η εξίσωση $f(x)=k$ έχει ακριβώς μία λύση στο $(0,1)$, διότι η $f$ είναι γνησίως αύξουσα από το $-\infty$ στο $0$.

\item Αν $k=0$, η εξίσωση έχει ακριβώς μία λύση:
\[
\ln x=0\Leftrightarrow {x=1}.
\]

\item Αν $0<k<\dfrac{1}{e}$, η εξίσωση έχει ακριβώς δύο λύσεις: μία στο $(1,e)$ και μία στο $(e,+\infty)$.

\item Αν $k=\dfrac{1}{e}$, η εξίσωση έχει ακριβώς μία λύση, την
\[
{x=e},
\]
στην οποία η $f$ παρουσιάζει το ολικό μέγιστό της.

\item Αν $k>\dfrac{1}{e}$, η εξίσωση δεν έχει λύσεις.
\end{itemize}

\item Για κάθε $x\in[1,e]$ είναι $\ln x\geq0$, οπότε $f(x)\geq0$. Άρα το ζητούμενο εμβαδόν είναι
\begin{align*}
E
&=\int_1^e\frac{\ln x}{x}\d x\\
&=\left.\frac{1}{2}\ln^2x\right|_{1}^{e}\\
&={\frac{1}{2}}.
\end{align*}
\end{erwthma}
\end{thema}
%=========== ΤΕΛΟΣ ΛΥΣΗΣ - 60ο ΘΕΜΑ Δ ===========
